% Revize výukové přípravy, výchozí soubor GitHub Work-MAT, 87. hodina.
% Zdroje: Realisticky.cz 020409–020410; návaznost na grafy v hodinách 70–71
% Úlohy jsou vlastní / úlohy výchozí přípravy, případně nově didakticky obměněné.
\documentclass[11pt,a4paper]{article}
\input{../../sablony/pripravy/priprava.sty}
\usepackage{tikz}
\nastavPripravu{1. ročník}{Rovnice a nerovnice}{87}{Rovnice s absolutní hodnotou}
\begin{document}
\blok{Návaznost: od grafu k rovnici}
Z hodin 70--71 známe tvar grafu $y=|x-a|$. Rovnice $|x-a|=r$ hledá průsečíky s vodorovnou přímkou $y=r$; graf vysvětlí, proč mohou být dva, jeden nebo žádný.
\begin{definiceBox}
\textbf{Rovnice s absolutní hodnotou} je rovnice, v níž se neznámá nachází uvnitř absolutní hodnoty. Řešení musí vyhovovat původní rovnici.
\end{definiceBox}
\blok{Dvě základní situace}
Pro konstantu $c$ používáme $|A(x)|=c$: pro $c<0$ není řešení, pro $c=0$ platí $A(x)=0$ a pro $c>0$ platí $A(x)=c$ nebo $A(x)=-c$.
Jestliže napravo stojí výraz $B(x)$, pak \textbf{nejdřív požadujeme $B(x)\ge0$}. Potom můžeme rozdělit definiční obor podle znaménka $A(x)$; každý kandidát musí patřit do příslušné části a projít zkouškou.
\blok{Společně řešené příklady}
\begin{prikladReseniBox}
Vyřešte $|x-3|=5$ a objasněte obě řešení jako vzdálenost na číselné ose.
\tcblower
\[\begin{aligned}
x-3&=5\quad\text{nebo}\quad x-3=-5 \\
x&=8\quad\text{nebo}\quad x=-2 \\
K&=\{-2;8\}
\end{aligned}\]
Obě čísla jsou od trojky vzdálena právě o pět jednotek.
\end{prikladReseniBox}
\begin{prikladReseniBox}
Vyřešte $|x-2|=x$ a ověřte nezápornost pravé strany.
\tcblower
\[\begin{aligned}
x&\ge0\quad\text{(nutná podmínka)} \\
x<2&:\quad 2-x=x\ \Rightarrow\ x=1 \\
x\ge2&:\quad x-2=x\ \Rightarrow\ -2=0 \\
K&=\{1\}
\end{aligned}\]
Zkouška: $|1-2|=1$.
\end{prikladReseniBox}
\begin{prikladReseniBox}
Vyřešte $|2x-1|=x+2$.
\tcblower
\[\begin{aligned}
x&\ge-2\quad\text{(pravá strana)} \\
x<\frac12&:\quad 1-2x=x+2\ \Rightarrow\ x=-\frac13 \\
x\ge\frac12&:\quad 2x-1=x+2\ \Rightarrow\ x=3 \\
K&=\left\{-\frac13;3\right\}
\end{aligned}\]
Oba kořeny splňují svoji intervalovou podmínku i původní rovnici.
\end{prikladReseniBox}
\blok{Samostatné procvičování}
\begin{zadaniReseniBox}
$|x+4|=2$.
\tcblower
\[\begin{aligned}
x+4&=\pm2 \\
x&=-6\quad\text{nebo}\quad x=-2 \\
K&=\{-6;-2\}
\end{aligned}\]

\end{zadaniReseniBox}
\begin{zadaniReseniBox}
$|3x-6|=0$.
\tcblower
\[\begin{aligned}
3x-6&=0 \\
x&=2 \\
K&=\{2\}
\end{aligned}\]

\end{zadaniReseniBox}
\begin{zadaniReseniBox}
$|x-1|=-2$.
\tcblower
\[\begin{aligned}
|x-1|&\ge0>-2 \\
K&=\varnothing
\end{aligned}\]

\end{zadaniReseniBox}
\begin{zadaniReseniBox}
$|x+3|=x+1$.
\tcblower
\[\begin{aligned}
x&\ge-1\quad\text{(nutná podmínka)} \\
x\ge-3&:\ x+3=x+1\ \Rightarrow\ 3=1 \\
x<-3&:\ -x-3=x+1\ \Rightarrow\ x=-2\not<-3 \\
K&=\varnothing
\end{aligned}\]

\end{zadaniReseniBox}
\blok{Volitelná rezerva}
\begin{zadaniReseniBox}
Srovnejte grafy $y=|x-2|$ a $y=x$. Kolik průsečíků mají a proč?
\tcblower
\[\begin{aligned}
|x-2|&=x \\
x&=1 \\
K&=\{1\}
\end{aligned}\]
Při $x<0$ má pravá strana zápornou hodnotu a průsečík vzniknout nemůže.
\end{zadaniReseniBox}
\blok{Domácí úkol}
\begin{zadaniReseniBox}[title={Domácí příprava}]
\begin{enumerate}[label=\alph*),itemsep=2pt]
\item $|2x-8|=4$.
\item $|x|=7$.
\item $|x+3|=x+1$.
\item $|x-4|=5$.
\item $|3x+6|=0$.
\item $|2x-1|=x+2$.
\end{enumerate}
\tcblower
\begin{enumerate}[label=\alph*),itemsep=1pt]
\item
\[\begin{aligned}
2x-8&=\pm4 \\
x&=2\ \text{nebo}\ x=6 \\
K&=\{2;6\}
\end{aligned}\]
\item
\[\begin{aligned}
x&=\pm7 \\
K&=\{-7;7\}
\end{aligned}\]
\item
\[\begin{aligned}
x&\ge-1\quad\text{nutná podmínka} \\
x+3&=x+1\quad\text{pro }x\ge-3\quad\text{-- spor} \\
-x-3&=x+1\ \Rightarrow\ x=-2\not<-3 \\
K&=\varnothing
\end{aligned}\]
\item
\[\begin{aligned}
x-4&=\pm5 \\
x&=-1\ \text{nebo}\ x=9 \\
K&=\{-1;9\}
\end{aligned}\]
\item
\[\begin{aligned}
3x+6&=0 \\
x&=-2 \\
K&=\{-2\}
\end{aligned}\]
\item
\[\begin{aligned}
x<\frac12&:\quad 1-2x=x+2\ \Rightarrow\ x=-\frac13 \\
x\ge\frac12&:\quad 2x-1=x+2\ \Rightarrow\ x=3 \\
K&=\left\{-\frac13;3\right\}
\end{aligned}\]\par Oba výsledky splňují $x+2\ge0$.
\end{enumerate}
\end{zadaniReseniBox}
\poznamka{45 minut: grafický význam 7, pravidla a dva vzory 14, třetí vzor 7, samostatná práce 13, kontrola 4; rezervní úloha je dobrovolná.}
\end{document}
